% Generated draft: requires mathematical and source review.
\documentclass[11pt]{article}
\usepackage[margin=25mm]{geometry}
\usepackage{fontspec}
\setmainfont{Latin Modern Roman}
\usepackage{amsmath,amssymb,mathtools,unicode-math}
\setmathfont{Latin Modern Math}
\usepackage{enumitem,needspace,xurl,hyperref}
\setlength{\parindent}{0pt}
\setlength{\parskip}{5pt}
\setlength{\emergencystretch}{3em}
\newcommand{\N}{\mathbb{N}}
\newcommand{\Z}{\mathbb{Z}}
\newcommand{\Q}{\mathbb{Q}}
\newcommand{\R}{\mathbb{R}}
\newcommand{\C}{\mathbb{C}}
\newcommand{\F}{\mathbb{F}}
\newcommand{\A}{\mathbb{A}}
\newcommand{\PP}{\mathbb{P}}
\newcommand{\E}{\mathbb{E}}
\newcommand{\eps}{\varepsilon}
\newcommand{\dd}{\,\mathrm{d}}
\begin{document}
\section*{UnsolvedMath ID: 30002304}
\textbf{Source problem code:} catalogue entry
\subsection{Definitions and mathematical statement}
Here \emph{MSC 2010} denotes the 2010 Mathematics Subject Classification. A classificatory scheme for mathematics is understood descriptively as a scheme organizing mathematical subjects into branches. The supplied record does not specify a particular scheme beyond its reference to MSC 2010, or define a criterion for a classification to be historically ``natural.''

\textbf{Research question.} Explain how the classificatory scheme for mathematics developed, and investigate whether what the supplied statement calls MSC 2010's pure/applied distinction reflects a natural historical organization of mathematical branches.

This is a historiographical research question rather than a specified mathematical proposition. The source record supplies no formal meaning for ``natural historical organization,'' no precise historical scope, and no explicit identification of the asserted pure/applied distinction within MSC 2010. That characterization is retained as part of the supplied question and is not independently established here. These unresolved qualifications prevent a precise theorem-style formulation without additional source interpretation.
\subsection{Short English statement}
Investigate the historical development of mathematical classification and whether the pure/applied distinction attributed to MSC 2010 in the supplied record reflects a historically natural organization of mathematical subjects. The intended historical scope and the meaning of ``natural'' require clarification.
\subsection{Sources}
\begin{itemize}
\item Local UnsolvedMath registry, \texttt{data/unsolvedmath/problems.json}, permanent ID \texttt{30002304}, source problem code \texttt{catalogue entry}, entitled ``Historical Foundations of Mathematical Classification.'' The record assigns the proposed year 2013 and names no individual proposer.
\item Supplied report citation: \emph{From ``Mixed'' to ``Applied'' Mathematics: Tracing an important dimension of mathematics and its history} (2013), Oberwolfach Reports, Report 12/2013, pp.\ 657--733, DOI 10.4171/owr/2013/12. \url{https://doi.org/10.4171/owr/2013/12}. The supplied background describes Oberwolfach Reports as documenting mathematical workshops and their open-problem discussions; no precise passage locating this question is supplied.
\item Supplied literature source: ``Oberwolfach Report 12/2013,'' \url{https://doi.org/10.4171/owr/2013/12}. The literature assessment dated \textbf{2026-08-21} labels the question open and describes the checked literature as providing competing historical analyses rather than a completed explanatory account. This is supplied dated curation, not an independently verified literature assessment; its time-sensitive status and cited claims require independent verification.
\item The registry labels the statement status ``exact'' and reports that its content was checked against the cited official report, with notation, context, and question form normalized without changing the source-level claim. These are supplied verification claims and have not been independently checked for this draft.
\end{itemize}
\end{document}
